%%%%%%%%%%%%%%%%
%%% lev 25 %%%
%%%%%%%%%%%%%%%%

\Chapter{25}

\Verse{1}
{
	\hP{וַיְדַבֵּר יְהֹוָה אֶל מֹשֶׁה בְּהַר סִינַי לֵאמֹר׃}
}
{
	\eP{And YHWH spoke to Moses in Mount Sinai, saying,}
}
{
}

\Verse{2}
{
	\hP{דַּבֵּר אֶל בְּנֵי יִשְׂרָאֵל וְאָמַרְתָּ אֲלֵהֶם כִּי תָבֹאוּ אֶל הָאָרֶץ אֲשֶׁר אֲנִי נֹתֵן לָכֶם וְשָׁבְתָה הָאָרֶץ שַׁבָּת לַיהֹוָה׃}
}
{
	\eP{
		“Speak to the children of Israel. And you shall say to
		them: When you will come to the land that I am giving to
		you, then the land shall have a Sabbath for YHWH.
	}
}
{
}

\Verse{3}
{
	\hP{שֵׁשׁ שָׁנִים תִּזְרַע שָׂדֶךָ וְשֵׁשׁ שָׁנִים תִּזְמֹר כַּרְמֶךָ וְאָסַפְתָּ אֶת תְּבוּאָתָהּ׃}
}
{
	\eP{
		Six years you shall seed your field, and six years you
		shall prune your vineyard, and you shall gather its
		produce.
	}
}
{
}

\Verse{4}
{
	\hP{וּבַשָּׁנָה הַשְּׁבִיעִת שַׁבַּת שַׁבָּתוֹן יִהְיֶה לָאָרֶץ שַׁבָּת לַיהֹוָה שָׂדְךָ לֹא תִזְרָע וְכַרְמְךָ לֹא תִזְמֹר׃}
}
{
	\eP{
		And in the seventh year the land shall have a Sabbath, a
		ceasing, a Sabbath for YHWH: you shall not seed your
		field, and you shall not prune your vineyard,
	}
}
{
}

\Verse{5}
{
	\hP{אֵת סְפִיחַ קְצִירְךָ לֹא תִקְצוֹר וְאֶת עִנְּבֵי נְזִירֶךָ לֹא תִבְצֹר שְׁנַת שַׁבָּתוֹן יִהְיֶה לָאָרֶץ׃}
}
{
	\eP{
		you shall not reap your harvest’s free growth, and you
		shall not cut off your untrimmed grapes. The land shall
		have a year of ceasing.
	}
}
{
%	\footnote{
%		free growth. Plants that sprout on their own after a harvest rather than
%		having been planted.
%	}
}

\Verse{6}
{
	\hP{וְהָיְתָה שַׁבַּת הָאָרֶץ לָכֶם לְאכְלָה לְךָ וּלְעַבְדְּךָ וְלַאֲמָתֶךָ וְלִשְׂכִירְךָ וּלְתוֹשָׁבְךָ הַגָּרִים עִמָּךְ׃}
}
{
	\eP{
		And the land’s Sabbath shall be yours for food: for you
		and for your servant and for your maid and for your
		employee and for your visitor who are residing with you
	}
}
{
}

\Verse{7}
{
	\hP{וְלִבְהֶמְתְּךָ וְלַחַיָּה אֲשֶׁר בְּאַרְצֶךָ תִּהְיֶה כל תְּבוּאָתָהּ לֶאֱכֹל׃}
}
{
	\eP{
		and for your domestic animal and for the wild animal that
		are in your land shall be all of its produce to eat.
	}
}
{
}

\Verse{8}
{
	\hP{וְסָפַרְתָּ לְךָ שֶׁבַע שַׁבְּתֹת שָׁנִים שֶׁבַע שָׁנִים שֶׁבַע פְּעָמִים וְהָיוּ לְךָ יְמֵי שֶׁבַע שַׁבְּתֹת הַשָּׁנִים תֵּשַׁע וְאַרְבָּעִים שָׁנָה׃}
}
{
	\eP{
		“And you shall count seven Sabbaths of years: seven years,
		seven times. And the days of the seven Sabbaths of years
		shall be for you forty-nine years.
	}
}
{
}

\Verse{9}
{
	\hP{וְהַעֲבַרְתָּ שׁוֹפַר תְּרוּעָה בַּחֹדֶשׁ הַשְּׁבִעִי בֶּעָשׂוֹר לַחֹדֶשׁ בְּיוֹם הַכִּפֻּרִים תַּעֲבִירוּ שׁוֹפָר בְּכל אַרְצְכֶם׃}
}
{
	\eP{
		And you shall pass a blasting horn in the seventh month on
		the tenth of the month: on the Day of Atonement you shall
		have the horn pass through all your land.
	}
}
{
}

\Verse{10}
{
	\hP{וְקִדַּשְׁתֶּם אֵת שְׁנַת הַחֲמִשִּׁים שָׁנָה וּקְרָאתֶם דְּרוֹר בָּאָרֶץ לְכל יֹשְׁבֶיהָ יוֹבֵל הִוא תִּהְיֶה לָכֶם וְשַׁבְתֶּם אִישׁ אֶל אֲחֻזָּתוֹ וְאִישׁ אֶל מִשְׁפַּחְתּוֹ תָּשֻׁבוּ׃}
}
{
	\eP{
		And you shall consecrate the year that makes fifty years
		and proclaim liberty in the land, to all its inhabitants.
		It shall be a jubilee for you. And you shall go back, each
		to his possession; and you shall go back, each to his
		family.
	}
}
{
%	\footnote{
%		possession. Meaning: each person’s ancestral property. Footnote 25:10. a
%		jubilee: you shall go back, each to his possession. In the law of the
%		jubilee, YHWH commands that every fifty years all property is to return
%		to the original owners. This appears to be an economic program designed
%		to prevent the feudal system, common in the rest of the ancient Near
%		East, from developing in Israel. That is, it functions to prevent the
%		establishment of a class of wealthy landowners at the top of the
%		economic scale and a mass of landless peasants at the bottom. Every
%		Israelite is to be apportioned some land (described in the books of
%		Numbers and Joshua), and the deity commands that in every fiftieth year
%		the system returns to where it started. If an Israelite has lost his
%		ancestral land as a result of debt or calamity, he regains ownership of
%		it in the jubilee year. Land is unalienable. Individuals can suffer
%		difficult times, but there is a divinely decreed limit to their loss,
%		and the nation as a whole can never degenerate into a two-tiered system
%		of the very rich and the very poor.
%	}
}

\Verse{11}
{
	\hP{יוֹבֵל הִוא שְׁנַת הַחֲמִשִּׁים שָׁנָה תִּהְיֶה לָכֶם לֹא תִזְרָעוּ וְלֹא תִקְצְרוּ אֶת סְפִיחֶיהָ וְלֹא תִבְצְרוּ אֶת נְזִרֶיהָ׃}
}
{
	\eP{
		It, the year that makes fifty years, shall be a jubilee
		for you. You shall not seed, and you shall not harvest its
		free growths, and you shall not cut off its untrimmed
		fruits,
	}
}
{
}

\Verse{12}
{
	\hP{כִּי יוֹבֵל הִוא קֹדֶשׁ תִּהְיֶה לָכֶם מִן הַשָּׂדֶה תֹּאכְלוּ אֶת תְּבוּאָתָהּ׃}
}
{
	\eP{
		because it is a jubilee; it shall be a holy thing to you.
		You shall eat its produce from the field.
	}
}
{
}

\Verse{13}
{
	\hP{בִּשְׁנַת הַיּוֹבֵל הַזֹּאת תָּשֻׁבוּ אִישׁ אֶל אֲחֻזָּתוֹ׃}
}
{
	\eP{
		In this jubilee year you shall go back, each to his
		possession.
	}
}
{
}

\Verse{14}
{
	\hP{וְכִי תִמְכְּרוּ מִמְכָּר לַעֲמִיתֶךָ אוֹ קָנֹה מִיַּד עֲמִיתֶךָ אַל תּוֹנוּ אִישׁ אֶת אָחִיו׃}
}
{
	\eP{
		And when you will make a sale to your fellow or buy from
		your fellow’s hand: each of you, do not persecute his
		brother.
	}
}
{
}

\Verse{15}
{
	\hP{בְּמִסְפַּר שָׁנִים אַחַר הַיּוֹבֵל תִּקְנֶה מֵאֵת עֲמִיתֶךָ בְּמִסְפַּר שְׁנֵי תְבוּאֹת יִמְכּר לָךְ׃}
}
{
	\eP{
		You shall buy from your fellow by the number of years
		after the jubilee. He shall sell to you by the number of
		years of yields—
	}
}
{
}

\Verse{16}
{
	\hP{לְפִי רֹב הַשָּׁנִים תַּרְבֶּה מִקְנָתוֹ וּלְפִי מְעֹט הַשָּׁנִים תַּמְעִיט מִקְנָתוֹ כִּי מִסְפַּר תְּבוּאֹת הוּא מֹכֵר לָךְ׃}
}
{
	\eP{
		corresponding to a greater amount of the years, you shall
		make its price greater; and corresponding to a smaller
		amount of the years, you shall make its price
		smaller—because it is a number of yields that he is
		selling to you.
	}
}
{
}

\Verse{17}
{
	\hP{וְלֹא תוֹנוּ אִישׁ אֶת עֲמִיתוֹ וְיָרֵאתָ מֵאֱלֹהֶיךָ כִּי אֲנִי יְהֹוָה אֱלֹהֵיכֶם׃}
}
{
	\eP{
		And you shall not persecute—each of you—his brother, but
		you shall fear your {\God}, because I am YHWH, your {\God}.
	}
}
{
}

\Verse{18}
{
	\hP{וַעֲשִׂיתֶם אֶת חֻקֹּתַי וְאֶת מִשְׁפָּטַי תִּשְׁמְרוּ וַעֲשִׂיתֶם אֹתָם וִישַׁבְתֶּם עַל הָאָרֶץ לָבֶטַח׃}
}
{
	\eP{
		“And you shall do my laws and observe my judgments and do
		them, so you will live on the land in security.
	}
}
{
}

\Verse{19}
{
	\hP{וְנָתְנָה הָאָרֶץ פִּרְיָהּ וַאֲכַלְתֶּם לָשֹׂבַע וִישַׁבְתֶּם לָבֶטַח עָלֶיהָ׃}
}
{
	\eP{
		And the land will give its fruit, and you will eat to the
		full, and you will live in security on it.
	}
}
{
}

\Verse{20}
{
	\hP{וְכִי תֹאמְרוּ מַה נֹּאכַל בַּשָּׁנָה הַשְּׁבִיעִת הֵן לֹא נִזְרָע וְלֹא נֶאֱסֹף אֶת תְּבוּאָתֵנוּ׃}
}
{
	\eP{
		And if you will say, ‘What shall we eat in the seventh
		year? Here, we won’t seed and won’t gather our produce,’
	}
}
{
}

\Verse{21}
{
	\hP{וְצִוִּיתִי אֶת בִּרְכָתִי לָכֶם בַּשָּׁנָה הַשִּׁשִּׁית וְעָשָׂת אֶת הַתְּבוּאָה לִשְׁלֹשׁ הַשָּׁנִים׃}
}
{
	\eP{
		I shall command my blessing for you in the sixth year, and
		it will make the produce for the three years.
	}
}
{
%	\footnote{
%		three years. The sixth year, the seventh year, and the eighth year (the
%		first year of the next seven-year cycle, until the first planting of
%		that year produces a crop).
%	}
}

\Verse{22}
{
	\hP{וּזְרַעְתֶּם אֵת הַשָּׁנָה הַשְּׁמִינִת וַאֲכַלְתֶּם מִן הַתְּבוּאָה יָשָׁן עַד הַשָּׁנָה הַתְּשִׁיעִת עַד בּוֹא תְּבוּאָתָהּ תֹּאכְלוּ יָשָׁן׃}
}
{
	\eP{
		And you will seed the eighth year and eat from the old
		produce until the ninth year; you will eat the old until
		its produce comes.
	}
}
{
}

\Verse{23}
{
	\hP{וְהָאָרֶץ לֹא תִמָּכֵר לִצְמִתֻת כִּי לִי הָאָרֶץ כִּי גֵרִים וְתוֹשָׁבִים אַתֶּם עִמָּדִי׃}
}
{
	\eP{
		But the land shall not be sold permanently, because the
		land is mine, because you are aliens and visitors with me!
	}
}
{
}

\Verse{24}
{
	\hP{וּבְכֹל אֶרֶץ אֲחֻזַּתְכֶם גְּאֻלָּה תִּתְּנוּ לָאָרֶץ׃}
}
{
	\eP{
		So in all the land in your possession you shall give a
		redemption for the land.
	}
}
{
}

\Verse{25}
{
	\hP{כִּי יָמוּךְ אָחִיךָ וּמָכַר מֵאֲחֻזָּתוֹ וּבָא גֹאֲלוֹ הַקָּרֹב אֵלָיו וְגָאַל אֵת מִמְכַּר אָחִיו׃}
}
{
	\eP{
		When your brother will be low so that he will sell any of his
		possession, then his redeemer who is the closest relative to him
		shall come and redeem his brother’s sale.
	}
}
{
%	\footnote{
%		25:25,35,39. when your brother will be low. Meaning: when your
%		fellow Israelite will be so poor that he has difficulty holding on
%		to his land. It has long been noted that the three times that this
%		phrase occurs in this chapter reflect three stages of increasing
%		hardship. In the first, the person has to sell his land. In the
%		second, he has both lost his land and is without money. And in the
%		third, he himself is sold (or sells himself) into servitude. And at
%		each stage, the man’s fellow Israelite is commanded to help him. If
%		he has to sell his land, one must redeem it for him and then give it
%		back to him in the jubilee year. If he is without money, one must
%		lend him money without any charge or interest. And if he is sold,
%		one must not treat him as a slave.
%
%		Here the units of law convey the specific requirements while the
%		arrangement conveys the basic principle, namely, that as one’s
%		brother’s need increases, so does one’s responsibility to help him.
%		Further, one must thus help one’s fellow as a matter of law. The
%		chapter never speaks of charity, nor does it appeal to one’s
%		feelings of compassion or generosity. An unfortunate Israelite need
%		not feel degraded to be poor nor ashamed to be pitied. Economic
%		suffering is rather treated as a reality of life, which one is
%		required by law to remedy. The poor man thus can know that his
%		brother is helping him because the system requires brothers to help
%		one another and that, if the shoe were on the other foot, he would
%		do the same for his brother. This is not to say that the text denies
%		or discourages feelings of compassion, but only that the fulfillment
%		of the law is not made dependent upon the presence or absence of
%		such feelings.
%	}
}

\Verse{26}
{
	\hP{וְאִישׁ כִּי לֹא יִהְיֶה לּוֹ גֹּאֵל וְהִשִּׂיגָה יָדוֹ וּמָצָא כְּדֵי גְאֻלָּתוֹ׃}
}
{
	\eP{
		And when a man will not have a redeemer, but his hand has
		come to attain it and he has found enough for his
		redemption,
	}
}
{
}

\Verse{27}
{
	\hP{וְחִשַּׁב אֶת שְׁנֵי מִמְכָּרוֹ וְהֵשִׁיב אֶת הָעֹדֵף לָאִישׁ אֲשֶׁר מָכַר לוֹ וְשָׁב לַאֲחֻזָּתוֹ׃}
}
{
	\eP{
		then he shall figure the years of his sale and pay back
		what is left over to the man to whom he sold it, and he
		shall go back to his possession.
	}
}
{
}

\Verse{28}
{
	\hP{וְאִם לֹא מָצְאָה יָדוֹ דֵּי הָשִׁיב לוֹ וְהָיָה מִמְכָּרוֹ בְּיַד הַקֹּנֶה אֹתוֹ עַד שְׁנַת הַיּוֹבֵל וְיָצָא בַּיֹּבֵל וְשָׁב לַאֲחֻזָּתוֹ׃}
}
{
	\eP{
		And if his hand has not found enough to pay it back to
		him, then his sale shall be in the hand of the one who
		bought it until the jubilee year, and it shall come out in
		the jubilee, and he shall go back to his possession.
	}
}
{
}

\Verse{29}
{
	\hP{וְאִישׁ כִּי יִמְכֹּר בֵּית מוֹשַׁב עִיר חוֹמָה וְהָיְתָה גְּאֻלָּתוֹ עַד תֹּם שְׁנַת מִמְכָּרוֹ יָמִים תִּהְיֶה גְאֻלָּתוֹ׃}
}
{
	\eP{
		“And when a man will sell a home in a walled city, then
		its redemption shall be until the end of the year of its
		sale. Its redemption shall be for days.
	}
}
{
%	\footnote{
%		for days. The word “days” is understood here to mean the rest of the
%		days of the year—up to one full year. His redemption period can amount
%		to days, not years.
%	}
}

\Verse{30}
{
	\hP{וְאִם לֹא יִגָּאֵל עַד מְלֹאת לוֹ שָׁנָה תְמִימָה וְקָם הַבַּיִת אֲשֶׁר בָּעִיר אֲשֶׁר [לוֹ] (לא) חֹמָה לַצְּמִיתֻת לַקֹּנֶה אֹתוֹ לְדֹרֹתָיו לֹא יֵצֵא בַּיֹּבֵל׃}
}
{
	\eP{
		And if it will not be redeemed by the completion of a full
		year for it, then the house that is in a city that has a
		wall shall be established permanently for the one who
		bought it, through his generations. It shall not come out
		in the jubilee.
	}
}
{
}

\Verse{31}
{
	\hP{וּבָתֵּי הַחֲצֵרִים אֲשֶׁר אֵין לָהֶם חֹמָה סָבִיב עַל שְׂדֵה הָאָרֶץ יֵחָשֵׁב גְּאֻלָּה תִּהְיֶה לּוֹ וּבַיֹּבֵל יֵצֵא׃}
}
{
	\eP{
		“As for houses of the villages that do not have a wall
		around them: it shall be figured along with the land’s
		field. It shall have redemption, and it shall come out in
		the jubilee.
	}
}
{
}

\Verse{32}
{
	\hP{וְעָרֵי הַלְוִיִּם בָּתֵּי עָרֵי אֲחֻזָּתָם גְּאֻלַּת עוֹלָם תִּהְיֶה לַלְוִיִּם׃}
}
{
	\eP{
		“As for the Levites’ cities, the houses in the cities of
		their possession: the Levites shall have eternal
		redemption.
	}
}
{
}

\Verse{33}
{
	\hP{וַאֲשֶׁר יִגְאַל מִן הַלְוִיִּם וְיָצָא מִמְכַּר בַּיִת וְעִיר אֲחֻזָּתוֹ בַּיֹּבֵל כִּי בָתֵּי עָרֵי הַלְוִיִּם הִוא אֲחֻזָּתָם בְּתוֹךְ בְּנֵי יִשְׂרָאֵל׃}
}
{
	\eP{
		And that which one from the Levites will redeem: it shall
		come out, a sale of a house and of a city of his
		possession, in the jubilee, because that—the houses of the
		Levites’ cities—is their possession among the children of
		Israel.
	}
}
{
}

\Verse{34}
{
	\hP{וּשְׂדֵה מִגְרַשׁ עָרֵיהֶם לֹא יִמָּכֵר כִּי אֲחֻזַּת עוֹלָם הוּא לָהֶם׃}
}
{
	\eP{
		But a field of their cities’ surrounding land shall not be
		sold, because it is an eternal possession for them.
	}
}
{
}

\Verse{35}
{
	\hP{וְכִי יָמוּךְ אָחִיךָ וּמָטָה יָדוֹ עִמָּךְ וְהֶחֱזַקְתָּ בּוֹ גֵּר וְתוֹשָׁב וָחַי עִמָּךְ׃}
}
{
	\eP{
		“And when your brother will be low so that his hand is
		slipping with you, then you shall take hold of him—an
		alien and a visitor—and he shall live with you.
	}
}
{
%	\footnote{
%		his hand is slipping. Meaning: he does not have enough food and money to
%		live. Footnote 25:35. an alien and a visitor. Being landless, he is now
%		like an alien or a visitor.
%	}
}

\Verse{36}
{
	\hP{אַל תִּקַּח מֵאִתּוֹ נֶשֶׁךְ וְתַרְבִּית וְיָרֵאתָ מֵאֱלֹהֶיךָ וְחֵי אָחִיךָ עִמָּךְ׃}
}
{
	\eP{
		Do not take interest or a charge from him; but you shall
		fear your {\God}, and your brother will live with you.
	}
}
{
}

\Verse{37}
{
	\hP{אֶת כַּסְפְּךָ לֹא תִתֵּן לוֹ בְּנֶשֶׁךְ וּבְמַרְבִּית לֹא תִתֵּן אכְלֶךָ׃}
}
{
	\eP{
		You shall not give him your money for interest, and you
		shall not give him your food for a charge.
	}
}
{
}

\Verse{38}
{
	\hP{אֲנִי יְהֹוָה אֱלֹהֵיכֶם אֲשֶׁר הוֹצֵאתִי אֶתְכֶם מֵאֶרֶץ מִצְרָיִם לָתֵת לָכֶם אֶת אֶרֶץ כְּנַעַן לִהְיוֹת לָכֶם לֵאלֹהִים׃}
}
{
	\eP{
		I am YHWH, your {\God}, who brought you out from the land of
		Egypt to give you the land of Canaan, to be {\God} to you.
	}
}
{
}

\Verse{39}
{
	\hP{וְכִי יָמוּךְ אָחִיךָ עִמָּךְ וְנִמְכַּר לָךְ לֹא תַעֲבֹד בּוֹ עֲבֹדַת עָבֶד׃}
}
{
	\eP{
		“And when your brother will be low with you so that he is
		sold to you, you shall not have him work a slave’s work.
	}
}
{
}

\Verse{40}
{
	\hP{כְּשָׂכִיר כְּתוֹשָׁב יִהְיֶה עִמָּךְ עַד שְׁנַת הַיֹּבֵל יַעֲבֹד עִמָּךְ׃}
}
{
	\eP{
		He shall be like an employee, like a visitor, with you. He
		shall work with you until the jubilee year.
	}
}
{
}

\Verse{41}
{
	\hP{וְיָצָא מֵעִמָּךְ הוּא וּבָנָיו עִמּוֹ וְשָׁב אֶל מִשְׁפַּחְתּוֹ וְאֶל אֲחֻזַּת אֲבֹתָיו יָשׁוּב׃}
}
{
	\eP{
		And he shall go out from you, he and his children with
		him, and go back to his family, and he shall go back to
		his fathers’ possession.
	}
}
{
}

\Verse{42}
{
	\hP{כִּי עֲבָדַי הֵם אֲשֶׁר הוֹצֵאתִי אֹתָם מֵאֶרֶץ מִצְרָיִם לֹא יִמָּכְרוּ מִמְכֶּרֶת עָבֶד׃}
}
{
	\eP{
		Because they are my servants, whom I brought out from the
		land of Egypt; they shall not be sold like the sale of a
		slave.
	}
}
{
%	\footnote{
%		servant, slave. The same Hebrew root (‘bd) is used for “servant” and
%		“slave.” It has a wider range of meaning in Hebrew than either word has
%		in English. Here it results in the use of its two different translations
%		in the very same sentence, but this still conveys the Torah’s meaning
%		more correctly than leveling the nuances of the Hebrew into a single
%		English word.
%	}
}

\Verse{43}
{
	\hP{לֹא תִרְדֶּה בוֹ בְּפָרֶךְ וְיָרֵאתָ מֵאֱלֹהֶיךָ׃}
}
{
	\eP{
		You shall not dominate him with harshness, but you shall
		fear your {\God}.
	}
}
{
%	\footnote{
%		dominate. “You shall not dominate him.” This is the term that is used
%		for human dominion over animals in Genesis 1. Israelites are explicitly
%		denied the kind of power over other human beings, even a slave, that
%		they might exert over animals. Footnote 25:43. with harshness. This is
%		the term that is used for the way the Egyptians made the Israelites work
%		(Exod 1:14). Eliminating slavery on earth has been a gradual process
%		rather than a revolution. The Israelites are permitted to have slaves
%		themselves, but not in the manner of the Egyptians. They cannot overwork
%		them. They cannot mistreat them (so an Israelite who knocks out his
%		slave’s tooth must set the slave free; Exod 21:2). The rules that are
%		imposed on treatment of slaves in the Torah involve a recognition of a
%		slave’s humanity and dignity, and they establish that a slave has
%		rights. This was a crucial first stage in the long process leading to
%		opposition to slavery altogether.
%	}
}

\Verse{44}
{
	\hP{וְעַבְדְּךָ וַאֲמָתְךָ אֲשֶׁר יִהְיוּ לָךְ מֵאֵת הַגּוֹיִם אֲשֶׁר סְבִיבֹתֵיכֶם מֵהֶם תִּקְנוּ עֶבֶד וְאָמָה׃}
}
{
	\eP{
		“And your slave and your maid that you will have: from the
		nations that are around you, you shall buy a slave or a
		maid from them.
	}
}
{
}

\Verse{45}
{
	\hP{וְגַם מִבְּנֵי הַתּוֹשָׁבִים הַגָּרִים עִמָּכֶם מֵהֶם תִּקְנוּ וּמִמִּשְׁפַּחְתָּם אֲשֶׁר עִמָּכֶם אֲשֶׁר הוֹלִידוּ בְּאַרְצְכֶם וְהָיוּ לָכֶם לַאֲחֻזָּה׃}
}
{
	\eP{
		And also from the children of the visitors who reside with
		you—you shall buy from them and from their families that
		are with you—to whom they gave birth in your land. And
		they shall become a possession for you.
	}
}
{
}

\Verse{46}
{
	\hP{וְהִתְנַחַלְתֶּם אֹתָם לִבְנֵיכֶם אַחֲרֵיכֶם לָרֶשֶׁת אֲחֻזָּה לְעֹלָם בָּהֶם תַּעֲבֹדוּ וּבְאַחֵיכֶם בְּנֵי יִשְׂרָאֵל אִישׁ בְּאָחִיו לֹא תִרְדֶּה בוֹ בְּפָרֶךְ׃}
}
{
	\eP{
		And you shall make them a legacy for your children after
		you to inherit as a possession. You may have them work
		forever, but as for your brothers, the children of
		Israel—a man toward his brother—you shall not dominate him
		with harshness.
	}
}
{
}

\Verse{47}
{
	\hP{וְכִי תַשִּׂיג יַד גֵּר וְתוֹשָׁב עִמָּךְ וּמָךְ אָחִיךָ עִמּוֹ וְנִמְכַּר לְגֵר תּוֹשָׁב עִמָּךְ אוֹ לְעֵקֶר מִשְׁפַּחַת גֵּר׃}
}
{
	\eP{
		“And if the hand of an alien and visitor with you will
		attain it, and your brother is low with him and is sold to
		the alien and visitor with you or to an offshoot of an
		alien’s family,
	}
}
{
%	\footnote{
%		attain it. That is, be able to provide the money to buy the poor
%		Israelite.
%	}
}

\Verse{48}
{
	\hP{אַחֲרֵי נִמְכַּר גְּאֻלָּה תִּהְיֶה לּוֹ אֶחָד מֵאֶחָיו יִגְאָלֶנּוּ׃}
}
{
	\eP{
		after he is sold he shall have redemption. One of his
		brothers shall redeem him,
	}
}
{
}

\Verse{49}
{
	\hP{אוֹ דֹדוֹ אוֹ בֶן דֹּדוֹ יִגְאָלֶנּוּ אוֹ מִשְּׁאֵר בְּשָׂרוֹ מִמִּשְׁפַּחְתּוֹ יִגְאָלֶנּוּ אוֹ הִשִּׂיגָה יָדוֹ וְנִגְאָל׃}
}
{
	\eP{
		or his uncle or his uncle’s son shall redeem him, or some
		close relative of his from his family shall redeem him; or
		if his hand can attain it then he can redeem himself.
	}
}
{
}

\Verse{50}
{
	\hP{וְחִשַּׁב עִם קֹנֵהוּ מִשְּׁנַת הִמָּכְרוֹ לוֹ עַד שְׁנַת הַיֹּבֵל וְהָיָה כֶּסֶף מִמְכָּרוֹ בְּמִסְפַּר שָׁנִים כִּימֵי שָׂכִיר יִהְיֶה עִמּוֹ׃}
}
{
	\eP{
		And he shall figure with the one who buys him: from the
		year he was sold to him to the jubilee year; and the price
		of his sale shall be by the number of years. It shall be
		with him like the days of an employee:
	}
}
{
}

\Verse{51}
{
	\hP{אִם עוֹד רַבּוֹת בַּשָּׁנִים לְפִיהֶן יָשִׁיב גְּאֻלָּתוֹ מִכֶּסֶף מִקְנָתוֹ׃}
}
{
	\eP{
		If there are still a great number of years, he shall pay
		back his redemption according to them from his purchase
		price.
	}
}
{
}

\Verse{52}
{
	\hP{וְאִם מְעַט נִשְׁאַר בַּשָּׁנִים עַד שְׁנַת הַיֹּבֵל וְחִשַּׁב לוֹ כְּפִי שָׁנָיו יָשִׁיב אֶת גְּאֻלָּתוֹ׃}
}
{
	\eP{
		And if a small amount of years remain until the jubilee
		year, then he shall figure it for him. He shall pay back
		his redemption according to his years.
	}
}
{
}

\Verse{53}
{
	\hP{כִּשְׂכִיר שָׁנָה בְּשָׁנָה יִהְיֶה עִמּוֹ לֹא יִרְדֶּנּוּ בְּפֶרֶךְ לְעֵינֶיךָ׃}
}
{
	\eP{
		He shall be with him like a year-by-year employee. He
		shall not dominate him with harshness before your eyes.
	}
}
{
}

\Verse{54}
{
	\hP{וְאִם לֹא יִגָּאֵל בְּאֵלֶּה וְיָצָא בִּשְׁנַת הַיֹּבֵל הוּא וּבָנָיו עִמּוֹ׃}
}
{
	\eP{
		And if he will not be redeemed by these, then he shall go
		out in the jubilee year, he and his children with him.
	}
}
{
}

\Verse{55}
{
	\hP{כִּי לִי בְנֵי יִשְׂרָאֵל עֲבָדִים עֲבָדַי הֵם אֲשֶׁר הוֹצֵאתִי אוֹתָם מֵאֶרֶץ מִצְרָיִם אֲנִי יְהֹוָה אֱלֹהֵיכֶם׃}
}
{
	\eP{
		Because the children of Israel are servants to me. They
		are my servants, whom I brought out from the land of
		Egypt. I am YHWH, your {\God}.
	}
}
{
}
